%\documentclass[spanish,a4paper,12pt,twosides]{book}
\documentclass[spanish,a4paper,12pt,oneside]{book}

\usepackage[T1]{fontenc}
\usepackage[latin1]{inputenc}
\usepackage[margin=1in]{geometry}
%\usepackage{paralist}
%\usepackage{graphicx}
\usepackage{url}
\usepackage{times}
\usepackage{babel}
\usepackage{listings}
\usepackage{verbatim}
\usepackage{float}
%\usepackage{color}
\usepackage{array}
\usepackage{fancyvrb}
\usepackage{relsize}
%\usepackage{eurosym}
%\usepackage[strict]{chngpage}
%\usepackage{qtree}
%\usepackage{longtable}
%\usepackage{cite}
\usepackage{titlesec}

\pagestyle{headings}

% Some subsvars
\newcommand{\TituloPFC}{Extracci�n de informaci�n de documentos XML mediante
Stanford Parser}
\newcommand{\TituloPFCWeb}{MIEX: Metadata and Information Extractor from small
XML documents}
\newcommand{\AutorPFC}{Ignacio Barrientos Arias}
\newcommand{\Escuela}{Escuela Universitaria de Ingenier�a T�cnica Inform�tica
de Gij�n}
\newcommand{\Universidad}{Universidad de Oviedo}
\newcommand{\Convocatoria}{SEPTIEMBRE 2007}

\newcommand{\sourcepath}
{/home/nacho/devel/pfc/svn/trunk/miex/src/es/uniovi/aic/miex/}
\newcommand{\trunkpath}
{/home/nacho/devel/pfc/svn/trunk/miex/}
\newcommand{\pfctrunkpath}
{/home/nacho/devel/pfc/svn/trunk/doc/pfc}
\newcommand{\codeversion}
{0.1.0}

% Some colour definitions
\definecolor{listinggray}{gray}{0.98}

% Cabeceras
\headheight 15pt

% Listing settings
\lstset
{
	basewidth=0.50em,
	backgroundcolor=\color{listinggray},
	basicstyle=\footnotesize\ttfamily,
	keywordstyle=\bfseries,
	stringstyle=\itshape,
	commentstyle=\itshape,
	showspaces=false,
	showtabs=false,
	showstringspaces=false,
	frame=trbl,
	extendedchars=true,
	%aboveskip=0.5cm,
	%belowskip=0.5cm,
	xleftmargin=0cm,
	xrightmargin=0cm,
	tabsize=1,
	numbers=left, 
	numberstyle=\small,
	stepnumber=5,
	numbersep=5pt
}

% Separation among paragraphs
\parskip  \medskipamount

% Titlesec
\titleformat{\chapter}[frame]
  {\fontfamily{phv}}
  {\filright\large\enspace \chaptertitlename \enspace \thechapter\enspace}
  {8pt}
  {\huge\bfseries\filcenter}
\titleformat{\section}
 {\normalfont\selectfont\large}
 {\thesection}{1em}{\bfseries}[\titlerule]
\titleformat{\section}
  {\fontfamily{phv}\selectfont\Large}
  {\thesection}{1em}{\bfseries}[\titlerule]
\titleformat{\subsection}
  {\fontfamily{phv}\selectfont\large}
  {\thesubsection}{1em}{\bfseries}

% Tweak to dynamically change the margin using a faked environment
\newenvironment{changemargin}[2]{
  \begin{list}{}{
    \setlength{\topsep}{0pt}
    \setlength{\leftmargin}{#1}
    \setlength{\rightmargin}{#2}
    \setlength{\listparindent}{\parindent}
    \setlength{\itemindent}{\parindent}
    \setlength{\parsep}{\parskip}
  }
\item[]}{\end{list}}


% Definici�n de macros para las portadas
\renewcommand{\maketitle}
{
  \begin{titlepage}

  % Portada
	\begin{center}

	\huge\textrm{\textbf{UNIVERSIDAD DE OVIEDO}}\\
	\vspace{1cm}

	\fbox{\includegraphics[width=3.5cm]{images/uniovi.png}}\\
	\vspace{1cm}

  \Large\textrm{\MakeUppercase{\Escuela}}\\
  \vspace{3cm}

  \Large\textrm{\textbf{PROYECTO FIN DE CARRERA}}\\
  \vspace{2cm}

  \Large\textrm{\textquotedblleft\MakeUppercase{\TituloPFC}\textquotedblright}\\
  \vspace{2cm}

  \Large\textrm{(C�DIGO FUENTE)}\\
  \vspace{1.5cm}

	\Large\textrm{AUTOR:\\ \AutorPFC}\\
	\vspace{0.5cm}
	
	\end{center}
	\begin{flushright}
		\large\textrm{\Convocatoria}
	\end{flushright}

        \pagestyle{empty}
	\cleardoublepage
    
  \end{titlepage}
    
  % Blank page before the main content
  \mbox{}
  \newpage
}


\title{\TituloPFC}
\author{\AutorPFC}
\date{Septiembre de 2007}

\begin{document}

% Arabic page numbers before the main content
\frontmatter

% Crear portadas
\maketitle

\chapter*{Aviso legal}

Copyleft 2007\hspace{2px}\textendash\hspace{2px}\AutorPFC{}

El c�digo fuente de todo el software escrito por el autor de este proyecto, al
que se hace referencia en numerosas ocasiones a lo largo de esta documentaci�n,
se encuentra bajo la licencia \emph{GNU General Public License (GPL)}, versi�n 2
(ap�ndice \ref{sec:license.gpl}). Las fuentes necesarias para generar la
documentaci�n se liberan bajo la misma licencia que el c�digo fuente del
programa.

\tableofcontents

\newpage

% Main stuff
\mainmatter

\chapter{C�digo fuente\label{sec:source}} 

La versi�n actual del c�digo (versi�n \codeversion) se incluye en este
ep�grafe. De cara a poner el software en producci�n se aconseja obtener la
�ltima versi�n de la p�gina web del proyecto\footnote{http://miex.sf.net}

\section*{Paquete Config}

\subsection*{ConfigCMDLine.java}

\lstinputlisting[language=Java]{\sourcepath/config/ConfigCMDLine.java}

\subsection*{ConfigFile.java}

\lstinputlisting[language=Java]{\sourcepath/config/ConfigFile.java}

\section*{Paquete Datastr}

\subsection*{ExtendedTaggedWord.java}
\lstinputlisting[language=Java]{\sourcepath/datastr/ExtendedTaggedWord.java}

\subsection*{MyCategories.java}
\lstinputlisting[language=Java]{\sourcepath/datastr/MyCategories.java}

\subsection*{MyCollection.java}
\lstinputlisting[language=Java]{\sourcepath/datastr/MyCollection.java}

\subsection*{MyDoc.java}
\lstinputlisting[language=Java]{\sourcepath/datastr/MyDoc.java}

\section*{Paquete filter}

\subsection*{GlobalFilter.java}
\lstinputlisting[language=Java]{\sourcepath/filter/GlobalFilter.java}

\subsection*{NumbersDetector.java}
\lstinputlisting[language=Java]{\sourcepath/filter/NumbersDetector.java}

\subsection*{Porter.java}
\lstinputlisting[language=Java]{\sourcepath/filter/Porter.java}

\subsection*{StopWordsDetector.java}
\lstinputlisting[language=Java]{\sourcepath/filter/StopWordsDetector.java}

\subsection*{UselessTagsDetector.java}
\lstinputlisting[language=Java]{\sourcepath/filter/UselessTagsDetector.java}

\section*{Paquete Input}

\subsection*{FieldsParser.java}
\lstinputlisting[language=Java]{\sourcepath/input/FieldsParser.java}

\subsection*{SAXCollectionUnmarshaller.java}
\lstinputlisting[language=Java]{\sourcepath/input/SAXCollectionUnmarshaller.java
}

\subsection*{XMLValidator.java}
\lstinputlisting[language=Java]{\sourcepath/input/XMLValidator.java}

\section*{Paquete Run}

\subsection*{Miex.java}
\lstinputlisting[language=Java]{\sourcepath/run/Miex.java}

\section*{Paquete Semantic}

\subsection*{Extractor.java}
\lstinputlisting[language=Java]{\sourcepath/semantic/Extractor.java}

\section*{Paquete Sql}

\subsection*{SQLHandler.java}
\lstinputlisting[language=Java]{\sourcepath/sql/SQLHandler.java}

\section*{Paquete Tools}

\subsection*{HexString.java}
\lstinputlisting[language=Java]{\sourcepath/tools/HexString.java}

\subsection*{MD5.java}
\lstinputlisting[language=Java]{\sourcepath/tools/MD5.java}

\section*{Otros}

\subsection*{Makefile (MIEX)}
\lstinputlisting{\trunkpath/Makefile}

\subsection*{Build.xml (MIEX)}
\lstinputlisting{\trunkpath/build.xml}

\subsection*{Makefile (Documentaci�n)}
\lstinputlisting{Makefile}


\end{document}
